\documentclass[a4paper,12pt]{article}

\usepackage{cmap}
\usepackage{mathtext}
\usepackage[T2A]{fontenc}
\usepackage[utf8]{inputenc}
\usepackage[english,russian]{babel}
\usepackage{indentfirst}
\usepackage{amsmath,amsfonts,amssymb,amsthm,mathtools}
\usepackage{icomma}
\usepackage{geometry}
\geometry{top=25mm}
\geometry{bottom=35mm}
\geometry{left=35mm}
\geometry{right=20mm}

\author{
Выполнил: Синев Д.\,Е. \\
Группа: М-15
}
\title{Лабораторная работа по \LaTeX}
\date{\today}

\begin{document}

\maketitle

\section*{Задание 5}
Пробелы внутри формул игнорируются.

Предположим, $x=y$, тогда \dots

Предположим, $$x=y,$$ тогда \dots

Предположим, $$x=y$$, тогда \dots

\section*{Задание 12}
\[
 x=1,...,100
\]
\[
 x=1,\dots,100
\]
\[
 1+\dots+100
\]
\[
 xy\dots z
\]
\[
 xy\cdots z
\]
\[
 x=1,\cdots,100
\]
\[
 1+\ldots+100
\]

\section*{Задание 19}
\[
( ),\ [ ],\ \{\},\ \langle\rangle
\]
\[
\bigl(\bigr),\ \Bigl(\Bigr),\ \biggl(\biggr),\ \Biggl(\Biggr)
\]
\[
\left(\frac{x+y}{x-z}\right)
\]

\section*{Задание 26}
\[
\bar x,\ \vec x,\ \hat x,\ \check x,\ \dot x,\ \ddot x,\ \dddot x,\ \tilde x
\]
\[
\overline{xyz},\ \overrightarrow{xyz},\ \widehat{xyz},\ \widetilde{xyz}
\]

\section*{Задание 33}
Пусть
\begin{equation}\label{eq:pq}
 p(x)=x^3+1.
\end{equation}
Из \eqref{eq:pq} следует, что \dots

\section*{Задание 40}
\[
\alpha x + \underbrace{\phi(y)}_{u}
+ \beta y + \underbrace{\psi(x)}_{v}
+ \Gamma(x)\times yx
\]
\[
\underbrace{\alpha x + \phi(y)}_{u}
+ \underbrace{\beta y + \psi(x)}_{v}
+ \underbrace{\Gamma(x)\times yx}_{w}
\]

\section*{Задание 47}
\[
\mu(n)=
\begin{cases}
0, & \text{если $n$ делится на квадрат простого числа};\\
(-1)^m, & \text{где $m$ — количество простых делителей $n$}.
\end{cases}
\]

\section*{Задание 54}
\[
 f^{(k)}(z_0)=\frac{k!}{2\pi i}\oint_{L_G}\frac{f(z)}{(z-z_0)^{k+1}}\,dz.
\]
\[
\text{Теорема Лиувилля: аналитическая и ограниченная на всем }\mathbb{C}\text{ функция является константой.}
\]
\[
\text{Основная теорема алгебры: каждый нетривиальный многочлен в }\mathbb{C}\text{ имеет корень.}
\]

\section*{Задание 7}
\[
\int xdx
\]
\[
\int x\,dx,\ \int x\:dx,\ \int x\;dx
\]
\[
 f(x)=2x,\quad g(x)=x-3
\]
\[
 x=x=x=x=x=x=x=x=x
\]

\section*{Задание 14}
\[
\forall \epsilon>0\dots |f(x)|<\epsilon
\]
\[
\forall \varepsilon>0\dots |f(x)|<\varepsilon
\]
\[
\let\epsilon\varepsilon\quad \forall \epsilon>0\dots |f(x)|<\epsilon
\]

\end{document}\documentclass[a4paper,12pt]{article}

\usepackage{cmap}
\usepackage{mathtext}
\usepackage[T2A]{fontenc}
\usepackage[utf8]{inputenc}
\usepackage[english,russian]{babel}
\usepackage{indentfirst}
\usepackage{amsmath,amsfonts,amssymb,amsthm,mathtools}
\usepackage{icomma}
\usepackage{geometry}
\geometry{top=25mm}
\geometry{bottom=35mm}
\geometry{left=35mm}
\geometry{right=20mm}

\author{
Выполнил: Синев Д.\,Е. \\
Группа: М-15
}
\title{Лабораторная работа по \LaTeX}
\date{\today}

\begin{document}

\maketitle

\section*{Задание 5}
Пробелы внутри формул игнорируются.

Предположим, $x=y$, тогда \dots

Предположим, $$x=y,$$ тогда \dots

Предположим, $$x=y$$, тогда \dots

\section*{Задание 12}
\[
 x=1,...,100
\]
\[
 x=1,\dots,100
\]
\[
 1+\dots+100
\]
\[
 xy\dots z
\]
\[
 xy\cdots z
\]
\[
 x=1,\cdots,100
\]
\[
 1+\ldots+100
\]

\section*{Задание 19}
\[
( ),\ [ ],\ \{\},\ \langle\rangle
\]
\[
\bigl(\bigr),\ \Bigl(\Bigr),\ \biggl(\biggr),\ \Biggl(\Biggr)
\]
\[
\left(\frac{x+y}{x-z}\right)
\]

\section*{Задание 26}
\[
\bar x,\ \vec x,\ \hat x,\ \check x,\ \dot x,\ \ddot x,\ \dddot x,\ \tilde x
\]
\[
\overline{xyz},\ \overrightarrow{xyz},\ \widehat{xyz},\ \widetilde{xyz}
\]

\section*{Задание 33}
Пусть
\begin{equation}\label{eq:pq}
 p(x)=x^3+1.
\end{equation}
Из \eqref{eq:pq} следует, что \dots

\section*{Задание 40}
\[
\alpha x + \underbrace{\phi(y)}_{u}
+ \beta y + \underbrace{\psi(x)}_{v}
+ \Gamma(x)\times yx
\]
\[
\underbrace{\alpha x + \phi(y)}_{u}
+ \underbrace{\beta y + \psi(x)}_{v}
+ \underbrace{\Gamma(x)\times yx}_{w}
\]

\section*{Задание 47}
\[
\mu(n)=
\begin{cases}
0, & \text{если $n$ делится на квадрат простого числа};\\
(-1)^m, & \text{где $m$ — количество простых делителей $n$}.
\end{cases}
\]

\section*{Задание 54}
\[
 f^{(k)}(z_0)=\frac{k!}{2\pi i}\oint_{L_G}\frac{f(z)}{(z-z_0)^{k+1}}\,dz.
\]
\[
\text{Теорема Лиувилля: аналитическая и ограниченная на всем }\mathbb{C}\text{ функция является константой.}
\]
\[
\text{Основная теорема алгебры: каждый нетривиальный многочлен в }\mathbb{C}\text{ имеет корень.}
\]

\section*{Задание 7}
\[
\int xdx
\]
\[
\int x\,dx,\ \int x\:dx,\ \int x\;dx
\]
\[
 f(x)=2x,\quad g(x)=x-3
\]
\[
 x=x=x=x=x=x=x=x=x
\]

\section*{Задание 14}
\[
\forall \epsilon>0\dots |f(x)|<\epsilon
\]
\[
\forall \varepsilon>0\dots |f(x)|<\varepsilon
\]
\[
\let\epsilon\varepsilon\quad \forall \epsilon>0\dots |f(x)|<\epsilon
\]

\end{document}